\section{From low energy to shifted signs}\label{sec:shifted-signs}

We prove Proposition~\ref{prop:adaptive}. Its spectral hypothesis will first
give a single normalized function with small derivatives in every coordinate. A lifted
Steiner symmetrization then produces a new domain from which one can take
one of two signed steps, even after a prescribed small shift. Repeating
this construction backwards will allow us to choose the signs forwards.
The construction adapts the common-density and prescribed-section arguments
of Guo, Fang and Lu~\cite[Proposition~2.2, Lemmas~3.3--3.4 and~4.1--4.3]{GuoFangLu2026}.
Quadratic-energy versions of this approach appear in
Bandeira~\cite{Bandeira2026} and Akbas and Sra
\cite[Lemmas~2.2--2.3]{AkbasSra2026}. Here we retain only coordinate Sobolev
energies and enough geometric slack for bounded adaptive shifts.
The coordinate-energy and shifted-step arguments needed here are proved
in full below.

Throughout this section,
\[
 \kappa=\frac1{100},\qquad \delta=\frac14,\qquad \ell=\frac52.
\]

\subsection{One function for all coordinate energies}

We first record a useful boundary convention. If a bounded open convex
domain $D\subset\R^m$ contains zero, then
\begin{equation}\label{eq:zero-extension}
 H_0^1(D)=\{f|_D:f\in H^1(\R^m),\ f=0\text{ a.e. outside }D\},
\end{equation}
where functions on the left are extended by zero. Indeed, for a function on
the right, $f_r(x)=f(x/r)$ converges to $f$ in $H^1(\R^m)$ as $r\uparrow1$.
For $r<1$ its support lies in the compact subset $r\overline D$ of $D$.
Mollification with a sufficiently small radius therefore approximates $f_r$
by functions in $C_c^\infty(D)$. Conversely, zero extension preserves the
$H^1$ norm on $C_c^\infty(D)$ and hence on its completion.

\begin{lemma}\label{lem:simultaneous-density}
Let $K\subset\R^m$ be a nonempty bounded open convex set symmetric about
zero. If
\[
 \lambda_Q(K)\le\kappa^2\tr Q
 \qquad\text{for every positive definite diagonal }Q,
\]
then there is a nonnegative $f\in H^1(\R^m)$, vanishing almost everywhere
outside $K$, such that
\[
 \int f^2=1,\qquad \int|\partial_i f|^2\le\kappa^2
 \quad(1\le i\le m).
\]
\end{lemma}

\begin{proof}
Let $\mathcal F$ be the nonnegative functions in $H^1(\R^m)$ that vanish
almost everywhere outside $K$ and have $L^2$ norm one, and put
$E_i(f)=\int|\partial_i f|^2$. Consider the set of simultaneous upper bounds
\[
 \mathcal A=\{u\in\R^m:E_i(f)\le u_i\text{ for every }i
                   \text{ for some }f\in\mathcal F\}.
\]
It is nonempty and is closed under addition of nonnegative vectors. It is
also convex. For $f,g\in\mathcal F$ and $0\le t\le1$, the function
$h=(tf^2+(1-t)g^2)^{1/2}$ belongs to $\mathcal F$ and satisfies
\[
 |\partial_i h|^2\le t|\partial_i f|^2+(1-t)|\partial_i g|^2
 \quad\text{a.e.}
\]
This follows from the Sobolev chain rule for the Euclidean norm and
Cauchy--Schwarz.

To see that $\mathcal A$ is closed, let $u^{(j)}\in\mathcal A$ converge to
$u$, and choose corresponding $f_j\in\mathcal F$. Their $H^1$ norms are
bounded. Since all supports lie in a fixed bounded set, weak $H^1$
compactness and compactness in $L^2$ give a subsequence converging weakly in
$H^1$ and strongly in $L^2$ to a member $f$ of $\mathcal F$. Lower
semicontinuity gives $E_i(f)\le u_i$ for every $i$.

If $c=\kappa^2(1,\ldots,1)$ did not belong to $\mathcal A$, strict
separation would give a nonzero $w\in\R^m$ with
\[
 w\cdot c<\inf_{u\in\mathcal A}w\cdot u.
\]
The upper-set property forces every $w_i$ to be nonnegative. For $\epsilon>0$
let $Q_\epsilon=\diag(w_i+\epsilon)$. By \eqref{eq:zero-extension}, taking
absolute values and normalizing does not change the Rayleigh infimum, so
\[
 \inf_{u\in\mathcal A}w\cdot u
 =\inf_{f\in\mathcal F}\sum_iw_iE_i(f)
 \le\lambda_{Q_\epsilon}(K)
 \le\kappa^2\sum_i(w_i+\epsilon).
\]
Letting $\epsilon\downarrow0$ contradicts separation. This argument also
covers zero components of $w$. Thus $c\in\mathcal A$, as required.
\end{proof}

\subsection{Convexity under Minkowski averages}

We will apply Jensen's inequality to the energies of horizontal sections.
The needed property is a diagonal version of the classical
Brascamp--Lieb convexity inequality for the first Dirichlet eigenvalue
\cite[Theorem~6.2]{BrascampLieb1976}; see also the statement in
Bryan, Clutterbuck and Rankin~\cite[Theorem~1.2]{BryanClutterbuckRankin2026}.
We give the proof from Gaussian log-concavity and the Brownian survival
rate, including the open-domain approximation used here.

\begin{lemma}\label{lem:minkowski}
Let $D_0,D_1\subset\R^m$ be nonempty bounded open convex sets symmetric
about zero, and let $Q$ be positive definite diagonal. For $0\le\alpha\le1$,
\[
 \lambda_Q((1-\alpha)D_0+\alpha D_1)
 \le(1-\alpha)\lambda_Q(D_0)+\alpha\lambda_Q(D_1).
\]
\end{lemma}

\begin{proof}
The endpoints are immediate, so suppose $0<\alpha<1$ and write
$D_\alpha=(1-\alpha)D_0+\alpha D_1$. Fix $0<\eta<1$, set $r=1-\eta$,
and put
\[
 A_i=r\overline{D_i},\qquad
 A_\alpha=(1-\alpha)A_0+\alpha A_1.
\]
These are compact convex sets, and $A_\alpha\subset D_\alpha$: indeed,
$A_\alpha\subset r\overline{D_\alpha}$, which is compactly contained in
$D_\alpha$.

Let $W$ start at zero with Brownian covariance rate $2Q$. On a finite grid
of positive times in $[0,T]$, its values form a Gaussian vector. Gaussian
log-concavity~\cite[Theorem~2]{Prekopa1973} shows that
\[
 \mathbb P\{W\text{ lies in }A_\alpha\text{ on the grid}\}
 \ge\prod_{i=0}^1
 \mathbb P\{W\text{ lies in }A_i\text{ on the grid}\}^{\alpha_i},
 \qquad (\alpha_0,\alpha_1)=(1-\alpha,\alpha).
\]
Here the Minkowski average of the two grid events is contained in the grid
event for $A_\alpha$. Increase the grids to a dense subset of $[0,T]$.
Closedness of the $A_i$ and continuity of Brownian paths allow passage to
the decreasing limits of these events. Since $rD_i\subset A_i$ and
$A_\alpha\subset D_\alpha$, it follows that
\[
 \mathbb P_0\{W([0,T])\subset D_\alpha\}
 \ge\prod_{i=0}^1
 \mathbb P_0\{W([0,T])\subset rD_i\}^{\alpha_i}.
\]
Lemma~\ref{lem:brownian-rate}, followed by the scaling of the Dirichlet
quotient, now gives
\[
 \lambda_Q(D_\alpha)
 \le r^{-2}\bigl((1-\alpha)\lambda_Q(D_0)+\alpha\lambda_Q(D_1)\bigr).
\]
Finally let $\eta\downarrow0$.
\end{proof}

\subsection{A section that permits shifted signed steps}

The next lemma is the geometric step in the signing construction. Its
conclusion holds for every allowed shift after the point in the new domain
has been chosen.

\begin{lemma}\label{lem:shift-section}
Let $K\subset\R^m$ be a nonempty bounded open convex set symmetric about
zero and satisfying $\lambda_Q(K)\le\kappa^2\tr Q$ for every positive
definite diagonal $Q$. For every coordinate vector $e$, there is another
domain $T_e(K)$ with the same properties such that, for all $y\in T_e(K)$
and $|q|\le\delta$, some $\varepsilon\in\{-1,1\}$ satisfies
\[
 y+(q+\varepsilon)e\in K.
\]
\end{lemma}

\begin{proof}
Lift $K$ to the bounded open convex set
\[
 B=\{(y,s)\in\R^m\times\R:y+se\in K,\ |s|<\ell\}.
\]
Let $P$ be its projection onto $\R^m$. For $y\in P$, its vertical fiber
$I_y=\{s:(y,s)\in B\}$ is a nonempty open interval. Write $L(y)=|I_y|$.
Convexity gives
\[
 (1-t)I_y+tI_z\subset I_{(1-t)y+tz},
 \qquad
 L((1-t)y+tz)\ge(1-t)L(y)+tL(z).
\]
Thus $L$ is positive, concave and continuous on the open convex set $P$.
The vertical Steiner symmetrization of $B$ is
\[
 B^*=\{(y,s):y\in P,\ |s|<L(y)/2\}.
\]
This set is bounded, open and convex. Symmetry of $K$ implies
$L(-y)=L(y)$, so $B^*$ is symmetric separately in $y$ and in $s$.
For $s\ge0$ put
\[
 D_s=\{y\in P:L(y)>2s\},\qquad
 b=\sup\{s\ge0:D_s\ne\varnothing\}.
\]
Then $0<b\le\ell$. Each $D_s$ with $0\le s<b$ is a nonempty bounded open
convex set symmetric about zero, and these domains decrease as $s$ increases.

The section $D_1$ is the target because each of its points permits the
required shifted signed step. Indeed, for any $y\in D_1$, write
$I_y=(a,b_0)$. Its length exceeds two. Choose
$c\in(a+1,b_0-1)$; then
\[
 [c-1,c+1]\subset I_y,\qquad |c|<\ell-1=\frac32.
\]
For $|q|\le\delta$, we have $|c-q|<7/4<2$. If $q\le c$, then $q+1$
lies in $[c-1,c+1]$; if $q\ge c$, then $q-1$ lies in that segment.
Thus one of $q+1,q-1$ lies in the original open fiber, giving membership
in $K$ itself. It remains to prove that $D_1$ is nonempty and has the same
spectral bounds as $K$.

Choose $f$ from Lemma~\ref{lem:simultaneous-density}, and define
\[
 \Psi(y,s)=(2\ell)^{-1/2}f(y+se)\mathbf1_{(-\ell,\ell)}(s).
\]
Translation invariance and Fubini give
\begin{equation}\label{eq:lifted-energy}
 \|\Psi\|_2=1,\qquad
 \int|\partial_{y_i}\Psi|^2=\int|\partial_i f|^2\le\kappa^2.
\end{equation}
Only derivatives in the horizontal variable $y$ will be used.

Rearrange each nonnegative fiber $\Psi(y,\cdot)$ into an even function
nonincreasing in $|s|$, denoted $\Psi^*(y,\cdot)$. A jointly measurable version
is obtained from the measurable superlevel lengths
$\mu(y,u)=|\{s:\Psi(y,s)>u\}|$ by replacing each superlevel set with its
centered interval of length $\mu(y,u)$. This preserves the $L^2$ norm.
Since $\Psi$ vanishes almost everywhere outside $B$, its rearrangement vanishes
almost everywhere outside $B^*$.

For nonnegative $a,b\in L^2(\R)$, the layer-cake formula gives
\[
 \int ab
 =\int_0^\infty\!\int_0^\infty
       |\{a>u\}\cap\{b>v\}|\,du\,dv
 \le\int a^*b^*.
\]
The intersection is bounded by the smaller of the two superlevel measures,
and centered intervals attain that bound. Since rearrangement preserves
the individual norms, $\|a^*-b^*\|_2\le\|a-b\|_2$. Applying this to the
fibers at $y$ and $y+he_i$, then integrating over $y$, bounds the horizontal
difference quotients of $\Psi^*$ by those of $\Psi$. Their weak limits yield
\begin{equation}\label{eq:rearranged-energy}
 \int|\partial_{y_i}\Psi^*|^2\le\int|\partial_i f|^2\le\kappa^2.
\end{equation}
Sobolev slicing and Fubini therefore imply that, for almost every $s$,
$\Psi^*(\cdot,s)\in H^1(\R^m)$ and vanishes almost everywhere outside
$D_{|s|}$. Each nonzero such slice belongs to $H_0^1(D_{|s|})$ by
\eqref{eq:zero-extension}.

We next show that the rearranged density puts enough mass at heights above
one. Its mean absolute height is
\begin{equation}\label{eq:mean-height}
 \begin{split}
 M&=\iint |s|\Psi^*(y,s)^2\,dy\,ds\\
 &=\frac1{8\ell}\int_{-\ell}^{\ell}\int_{-\ell}^{\ell}
      \int_{\R^m}\min\{f(y+re)^2,f(y+r'e)^2\}\,dy\,dr\,dr'.
 \end{split}
\end{equation}
Indeed, a centered interval of length $u$ has integral of $|s|$ equal to
$u^2/4$. Apply this to the superlevel intervals of $\Psi^{*2}$ and use Tonelli;
the factor $(2\ell)^{-1}$ in $\Psi^2$ gives the displayed normalization.
For nonnegative translates $u,v$ of $f$, both of norm one,
\[
 \int\min(u^2,v^2)
 =1-\tfrac12\|u^2-v^2\|_1
 \ge1-\|u-v\|_2.
\]
Moreover,
\[
 \|f(\cdot+re)-f(\cdot+r'e)\|_2
 \le |r-r'|\,\|\partial_e f\|_2\le2\ell\kappa.
\]
It follows from \eqref{eq:mean-height} that
\begin{equation}\label{eq:height-lower}
 M\ge\frac\ell2(1-2\ell\kappa)=\frac{19}{16}>1.
\end{equation}
Under the probability density $\Psi^{*2}(y,s)$, we have $|s|<b$ almost surely;
hence $M<b$. In particular, $D_1$ is nonempty.

It remains to pass the derivative bounds to this particular section.
Fix a positive definite diagonal $Q$ and put $h_Q(s)=\lambda_Q(D_s)$ for
$0\le s<b$. This function is finite and nondecreasing. Convexity of $B^*$,
domain monotonicity and Lemma~\ref{lem:minkowski} give
\[
 \begin{split}
 h_Q((1-t)s+tr)
 &\le\lambda_Q((1-t)D_s+tD_r)\\
 &\le(1-t)h_Q(s)+t h_Q(r),
 \end{split}
\]
so $h_Q$ is convex. Give $s$ the probability density
$\rho(s)=\int \Psi^*(y,s)^2\,dy$. Testing the Rayleigh quotient with each
nonzero slice and then integrating, we obtain
\[
 \int h_Q(|s|)\rho(s)\,ds
 \le\iint\langle Q\nabla_y\Psi^*,\nabla_y\Psi^*\rangle
 \le\kappa^2\tr Q.
\]
This also proves the integrability needed for Jensen's inequality. Since
$1<M<b$, monotonicity and Jensen give
\[
 \lambda_Q(D_1)=h_Q(1)\le h_Q(M)
 \le\int h_Q(|s|)\rho(s)\,ds\le\kappa^2\tr Q.
\]
The construction of $D_1$ is independent of $Q$, so every required spectral
bound is preserved. Taking $T_e(K)=D_1$ proves the lemma.
\end{proof}

\subsection{Choosing the signs}

\begin{proof}[Proof of Proposition~\ref{prop:adaptive}]
Write the prescribed coordinate directions as $e_{j_1},\ldots,e_{j_s}$.
Construct domains backwards by
\[
 K_s=K,\qquad K_{i-1}=T_{e_{j_i}}(K_i)
 \quad(i=s,s-1,\ldots,1).
\]
Lemma~\ref{lem:shift-section} preserves all domain and spectral hypotheses.
Each $K_i$ therefore contains zero, and we may start with $x^{(0)}=0\in K_0$.
Suppose $x^{(i-1)}\in K_{i-1}$, and set
\[
 q=q_i(x^{(i-1)}_{j_1},\ldots,x^{(i-1)}_{j_{i-1}}).
\]
Since $q\in[-\delta,\delta]$, the lemma supplies $\varepsilon_i\in\{-1,1\}$ with
\[
 x^{(i)}=x^{(i-1)}+(q+\varepsilon_i)e_{j_i}\in K_i.
\]
Induction gives $x^{(s)}\in K_s=K$ and the required signs. Distinctness of the
coordinate directions ensures that later steps do not alter the preceding
coefficient coordinates on which a shift depends. The backwards domains may
depend on the entire input sequence, as allowed in the proposition.
\end{proof}
